% =====================================================================
% Mechanical Design + Computer Vision  --  TARGET: ~600 words
% Marking weight: 10 marks (CAD/3D-printing). Vision content also
% supports the Software & Architecture rubric rows.
% =====================================================================
\section{Mechanical Design and Computer-Vision Subsystem}
\label{sec:cad}

\subsection{Mechanical design and 3D printing}
\guide{
For the 1st+ CAD rubric --- ``exceptional CAD design, optimised for
stability, usability, real-world application'' --- you need to show:
\begin{itemize}
  \item That the CAD is yours (or substantially modified), not
        printed unchanged from a public STL. Cite the original
        source~\cite{eezybotarm} but list the modifications you
        made.
  \item Print parameters: material, layer height, infill type \&
        \%, supports.
  \item Material justification: why PLA over PETG / ABS / nylon.
  \item DFM iterations: what failed in the first print, what you
        changed for the second.
\end{itemize}
A SolidWorks render and a photo of the printed parts mid-assembly
both pick up easy marks.
}

\placeholder{[\ldots one paragraph on the design intent and the
modifications made over the EEZYbotARM~\cite{eezybotarm} reference:
e.g.\ redesigned base plate, custom potentiometer mounts, wider
gripper opening, integrated cable channels\ldots]}

% TODO FIGURE: SolidWorks isometric render
\begin{figure}[t]
  \centering
  \fbox{\parbox{0.9\columnwidth}{\centering\vspace{4pt}
  \emph{[Isometric CAD render of the arm assembly. Export from
  SolidWorks as PDF or PNG.]}\vspace{4pt}}}
  \caption{Isometric CAD render (placeholder).}
  \label{fig:cadrender}
\end{figure}

\placeholder{[\ldots one paragraph on print parameters: PLA, layer
height, infill, perimeters, supports. Explain material choice and
mention the DFM tweaks made between iterations\ldots]}

% TODO PHOTO: printed parts mid-assembly
\begin{figure}[t]
  \centering
  \fbox{\parbox{0.9\columnwidth}{\centering\vspace{4pt}
  \emph{[Photograph of the 3D-printed parts laid out on the bench
  before assembly, or a half-assembled shot showing internal
  structure.]}\vspace{4pt}}}
  \caption{Printed parts (photo placeholder).}
  \label{fig:printedparts}
\end{figure}

\subsection{Computer-vision subsystem}
\label{sec:vision}
\guide{
Two-iteration narrative scores well in the 1st+ band because it
demonstrates engineering judgement, not just first-attempt success.
Structure:
\begin{enumerate}
  \item ESP32 first iteration: what worked (on-device HSV, UART link
        to Mega, no Bluetooth contention), what failed (lighting
        sensitivity, green/yellow boundary mis-classification).
  \item Decision to migrate to a Pi + camera module: explain the
        reasoning (more compute, larger frames, floating-point HSV,
        classical CV pipeline).
  \item Why the Mega side didn't change: same one-line ASCII
        protocol, same UART pin --- isolation between vision and
        motor control is preserved.
\end{enumerate}
This sub-section also evidences the Pi-vs-Arduino position from
Section~\ref{sec:literature}: Pi is fine \emph{as a co-processor}
but never owns real-time motor control.
}

\subsubsection{First iteration --- ESP32 WROVER CAM}
\placeholder{[\ldots one paragraph on the ESP32 hardware (Freenove
WROVER CAM, OV2640 sensor), the on-device HSV pipeline (QQVGA
RGB565 $\rightarrow$ integer HSV $\rightarrow$ four colour bins with
a 4\,\% noise floor), and the UART link
(GPIO~17 $\rightarrow$ Mega pin~19, no level shift required).
Mention initial successful classification.\ldots]}

% TODO PHOTO: ESP32 + camera mounted on bench during testing
\begin{figure}[t]
  \centering
  \fbox{\parbox{0.9\columnwidth}{\centering\vspace{4pt}
  \emph{[Photograph of the ESP32 WROVER CAM during the first
  iteration --- ideally with the laptop showing live colour
  classification.]}\vspace{4pt}}}
  \caption{ESP32 WROVER CAM during the first vision iteration
  (photo placeholder).}
  \label{fig:esp32cam}
\end{figure}

\placeholder{[\ldots one short paragraph on what failed: ambient
lighting sensitivity, occasional false classifications at the
green/yellow boundary, the 4\,\% noise-floor proving insufficient
under fluorescent lighting\ldots]}

\subsubsection{Second iteration --- Raspberry Pi camera module}
\placeholder{[\ldots one paragraph on the Pi pipeline: CSI camera
module, larger frame size, floating-point HSV conversion, classical
CV (Gaussian blur $\rightarrow$ HSV thresholding $\rightarrow$
contour-area filtering), and that the same ASCII protocol on the
same Mega UART pin meant zero firmware changes on the arm
side\ldots]}

% TODO PHOTO: Pi + camera module mounted on arm
\begin{figure}[t]
  \centering
  \fbox{\parbox{0.9\columnwidth}{\centering\vspace{4pt}
  \emph{[Photograph of the Raspberry Pi + camera module during the
  second iteration.]}\vspace{4pt}}}
  \caption{Raspberry Pi + camera module (photo placeholder).}
  \label{fig:picam}
\end{figure}

\placeholder{[\ldots closing paragraph reinforcing that all real-time
motor control remains on the Mega; the Pi is a vision co-processor
only. This preserves the embedded-systems character of the
project.\ldots]}

\corehit{Brief Objective~3 and the CAD/3D-printing rubric: original
CAD modifications, FDM print parameters justified, fabrication
performed in-house.}

\advhit{[\ldots e.g.\ two iterations of an embedded computer-vision
pipeline (on-device HSV on ESP32; classical CV on Pi); shared
single-protocol UART link so the arm-side firmware is unchanged
between iterations; multi-MCU heterogeneous architecture with clean
separation of real-time and non-real-time concerns.\ldots]}
